\scanpage{31}
$\overline{\mathcal V}=$ catégorie des var. top. \inserted{à bord} avec immersions ouvertes

$\mathcal V^{\unsure{*}}=$ catégorie des variétés topologiques sans bord

\diagram{\begin{tikzcd}\partial:\overline{\mathcal V}\arrow[r]&\mathcal V\end{tikzcd}} foncteur bord topologique,
\margin{\diagram{\begin{tikzcd}\unsure{X|\widehat X}\arrow[r]&\mathrm{int}(X)\inserted{=X-\partial X}\end{tikzcd}}\quad$\unsure{/\,=|:}$\quad\diagram{\begin{tikzcd}\overline{\mathcal V}\arrow[r]&\mathcal V\end{tikzcd}}\\foncteur ``\unsure{pts} intérieurs}

$\overline{\mathcal V}_{\le n},\mathcal V_{\le n},\overline{\mathcal V}_n,\mathcal V_n$ ss-catégories des variétés de dim $\le n$ (\unsure{resp.} purement de dim $n$)

Une \uline{théorie de l'orientation} des variétés topologiques consiste en

(i) Un foncteur
\[
\Omega:\quad\diagram{\begin{tikzcd}\mathcal V^0\arrow[r]&(\mathrm{Ens})\end{tikzcd}}
\]
\margin{On note $\Omega_n=\Omega\mid\mathcal V_n$\\$\Omega(X)=$ ensemble des orientations de $X$}

(ii) Une \deleted{\unsure{application|morphisme}} élément \unsure{de}
\[
\omega^+\in\Omega(e)\qquad(e\text{ la variété ponctuelle})
\]

(iii) \deleted{\unsure{Donnant}} Un hom de foncteurs sur $\overline{\mathcal V}^{0}$
\diagram{\begin{tikzcd}\Omega(\mathrm{int}\,X)\arrow[r]&\Omega(\partial X),\end{tikzcd}}

Ces données satisfaisant aux conditions suivantes

1°) $\Omega$ est un \uline{faisceau} [i.e $\forall X\in\mathrm{Ob}\mathcal V$, la \unsure{restrict.} \inserted{$\Omega_X$} de $\Omega$ à la catégorie des ouverts de $X$ est un faisceau (faisceau d'orientation de $X$)]
\margin{nature locale}
\deleted{\unsure{pour que} $\Omega_X\to X$ \unsure{soit} \illegible{1} (\unsure{tt ouvert} \illegible{3} \unsure{orienté})}

2°) $\mathrm{card}(\Omega(e))=2$. [L'élément de $\Omega(e)$ distinct de $\omega^+$ est noté $\omega^-$. On définit une bijection
\[
\Omega(e)\simeq\{+1,-1\}\subset\mathbb Z\quad\text{par}\quad\omega^+\mapsto+1,\ \omega^-\mapsto-1:
\]
multiplicité (égale à $+1$ ou $-1$) d'un point orienté]
\margin{hyp. de \unsure{démarrage|démarr.} ou récurrence (normalisation)}

\inserted{Catégorie $\vec{\mathcal V}^{(\unsure{\Omega|\omega})}$ des variétés top. ``orientées'' (i.e $\Omega$-orientées)}

\deleted{\unsure{3}} Pour tte $\vec X\in\mathrm{Ob}\vec{\mathcal V}$ et $Y$ sous-variété de codim 1, le faisceau $\rho_{Y/X}^{\unsure{(X'',\omega)|(X',\omega)}}$ des ``rives'' de $Y$ dans $X$ (qui est un rev. étale d'ordre 2 de $X$) est envoyé, grâce à la donnée (iii), dans le faisceau d'orientation $\Omega_Y$. On trouve donc, \deleted{pour tout $X\in$}
\diagram{\begin{tikzcd}\Omega(X)\arrow[r]&\mathrm{Hom}(\rho_{Y/X},\Omega_Y),\end{tikzcd}}
et par faisceautisation
\diagram{\begin{tikzcd}\Omega_X\mid Y\arrow[r]&\uline{\mathrm{Hom}}(\unsure{\rho|\varphi}_{Y/X},\Omega_Y)\end{tikzcd}}
ou encore
\diagram{\begin{tikzcd}\Omega_X\times\rho_{Y/X}\arrow[r]&\Omega_Y\end{tikzcd}}
\margin{30}
